\begin{frame}{Ergebnisvideo und Blur-Qualitaet}
\begin{itemize}
    \item Ovaler Gaussian Blur wirkt natuerlich und deckt Gesichter weich ab.
    \item Pixelate ist sichtbarer und kann bei kleinen Gesichtern haerter wirken.
    \item Die Blur-Qualitaet haengt direkt am Detektions-Recall.
    \item Smoke-Videos zeigen die Pipeline, ersetzen aber keine finale Datenschutzpruefung.
\end{itemize}
\vspace{0.6em}
\begin{mynote}[Qualitaetscheck]
Technisch gute Metriken reichen nicht: Am Ende muss das Video visuell auf verpasste Gesichter geprueft werden.
\end{mynote}
\end{frame}
